\section{标准形式与线性规划}
\label{sec:09-Convex-Optimization/04-Convex-Problems/01}

手上是一份排产表：七条产线、三十几种物料、每周重排一次。业务方把规则口述给你——某些物料只能上指定产线，每条线每天有工时上限，总产量不得低于订单量，成本压到最低。你写了个脚本，随机初始化加人工微调，跑一夜出一个方案，交付的时候只能说``这是我搜到的比较好的排法''。

半年后，另一个人把同一件事写成了四行：一个线性目标、一组线性不等式、一组线性等式、变量非负。丢进求解器，两秒返回，附带一个数字——这个方案与理论最优值的差距是零。差别不在谁的算法更聪明。差别在于问题被写成了什么形状：一种形状只能靠搜，另一种形状可以直接交给一个别人已经打磨了几十年的程序，并且拿回一份可以核验的最优性说明。

这个``形状''就是凸优化的\textbf{标准形式}：

\[
\begin{aligned}
\text{minimize}\quad &f_0(x)\\
\text{subject to}\quad
&f_i(x)\le0,\quad i=1,\ldots,m,\\
&Ax=b,
\end{aligned}
\]

其中 \(f_0,f_1,\ldots,f_m\) 都是凸函数，等式约束仿射。三条要求看着像格式规定，其实每一条都在买一样具体的东西。

\subsection{三条要求各买下一条几何性质}

不等式那一条买的是可行域的凸性。凸函数的次水平集 \(\set{x:f_i(x)\le0}\) 必然凸，若干个凸集取交仍然凸，所以 \(m\) 条不等式叠在一起不会把可行域割出凹口。方向反过来写成 \(f_i(x)\ge0\)，得到的是超水平集，除非 \(f_i\) 恰好是凹的，否则集合会朝内塌陷。一个负号写错，问题的全局性质就换了一副面孔。

等式那一条买的是同一件事，只是更苛刻。\(h(x)=0\) 通常给出一张曲面，而曲面上任取两点，连线一般要脱离曲面。\(x_1^2+x_2^2-1=0\) 就是现成的例子：左边这个函数凸得无可挑剔，它的零集却是单位圆周，圆周上两点的连线穿过圆盘内部。只有 \(h\) 仿射时，零集才是超平面，线段才留在里面。

目标那一条买的是\hyperref[sec:09-Convex-Optimization/01-Optimization-Foundations/02]{``局部更好为何还不够''}里那条定理的前提：凸目标加凸可行域，局部最优即全局最优。三条合起来的效果是，任何一个满足一阶最优性条件的点都可以被当作答案交出去，不必再担心它是不是某个浅坑。

所以判定一个问题凸不凸，不能只盯着目标函数。要逐项核对：最小化的是凸函数（或最大化的是凹函数），每条不等式以``凸函数 \(\le0\)''的方向出现，等式约束仿射，并且所有函数共同的隐式定义域也是凸的——最后这条最容易漏，\(\log\) 和矩阵求逆都会悄悄带进定义域限制。

\subsection{线性目标把最优解逼到顶点}

把标准形式里所有函数都退化成仿射的，得到\textbf{线性规划}：

\[
\begin{aligned}
\text{minimize}\quad &c^{\trans}x\\
\text{subject to}\quad
&Gx\le h,\\
&Ax=b.
\end{aligned}
\]

可行域是有限多个半空间与超平面的交，也就是一个\textbf{多面体}。目标的等值集 \(\set{x:c^{\trans}x=k}\) 是一族平行超平面，改变 \(k\) 就是把这族平面整体平移。求解在几何上只有一个动作：沿 \(-c\) 方向推这族平面，推到与多面体最后一次接触为止。

\begin{center}
\begin{tikzpicture}[x=1cm,y=1cm,font=\small,>=stealth]
  \node[font=\footnotesize,black!70] at (2.3,4.35) {等值线斜着推：最后落在一个顶点};
  \fill[black!8] (0,0) -- (3,0) -- (3.6,2) -- (2,3.6) -- (0,2.4) -- cycle;
  \begin{scope}
    \clip (-0.2,-0.2) rectangle (4.9,3.95);
    \draw[black!30] (0,2) -- (1,0);
    \draw[black!30] (0.55,3.9) -- (2.5,0);
    \draw[black!30] (1.55,3.9) -- (3.5,0);
    \draw[black!60,dashed] (2.65,3.9) -- (4.6,0);
  \end{scope}
  \draw[black!75,thick] (0,0) -- (3,0) -- (3.6,2) -- (2,3.6) -- (0,2.4) -- cycle;
  \fill[black] (3.6,2) circle (2.2pt);
  \node[font=\footnotesize] at (4.2,2.35) {\(x^\star\)};
  \draw[->,black!65,thick] (0.35,0.35) -- (1.15,0.75);
  \node[font=\footnotesize,black!65] at (1.5,0.7) {\(-c\)};
  \begin{scope}[shift={(6.4,0)}]
    \node[font=\footnotesize,black!70] at (2.3,4.35) {等值线与一条边平行：最优解不唯一};
    \fill[black!8] (0,0) -- (3,0) -- (3.6,2) -- (2,3.6) -- (0,2.4) -- cycle;
    \begin{scope}
      \clip (-0.2,-0.2) rectangle (4.9,3.95);
      \draw[black!30] (0,2) -- (2,0);
      \draw[black!30] (0.1,3.9) -- (4,0);
      \draw[black!60,dashed] (1.7,3.9) -- (4.9,0.7);
    \end{scope}
    \draw[black!75,thick] (0,0) -- (3,0) -- (3.6,2) -- (2,3.6) -- (0,2.4) -- cycle;
    \draw[black,line width=1.6pt] (3.6,2) -- (2,3.6);
    \fill[black] (3.6,2) circle (2.2pt);
    \fill[black] (2,3.6) circle (2.2pt);
    \node[font=\footnotesize] at (3.9,3.3) {这条边整段最优};
    \draw[->,black!65,thick] (0.35,0.35) -- (1.05,1.05);
    \node[font=\footnotesize,black!65] at (1.4,1.0) {\(-c\)};
  \end{scope}
\end{tikzpicture}

\emph{左：等值线是一族平行直线，沿 \(-c\) 推到最后一次接触，落点是一个顶点。右：等值线恰好与某条边平行时，这条边上每一点的目标值都相同。}
\end{center}

图里那族灰线越往右下走目标值越小，虚线是最后还碰得到可行域的那一条，它只压住一个角。这不是画图时挑出来的巧合。

\begin{proposition}
设线性规划可行、最优值有限，且可行多面体不含整条直线（变量非负这类约束已经保证了这一点），则最优值必在某个顶点取到。
\end{proposition}

骨架三步。取一个最优点 \(x^\star\)，若它不是顶点，则存在方向 \(d\ne0\) 使 \(x^\star\pm td\) 在足够小的 \(t>0\) 下都可行。目标沿这条直线是 \(c^{\trans}x^\star\pm t\,c^{\trans}d\)，若 \(c^{\trans}d\ne0\) 就有一侧严格更低，与最优矛盾；故 \(c^{\trans}d=0\)，沿 \(d\) 走目标纹丝不动。于是一直走下去，直到某条原本松弛的约束变紧为止——可行域不含直线保证这件事一定发生。新的点仍然最优，紧约束却多了一条。重复有限次，紧约束凑够 \(n\) 条独立的，就到了顶点。

条件对一遍账。第一步用的是目标线性：沿任何直线走，目标要么单调要么恒定，没有第三种可能，换成一般凸目标这一步立刻失效。第二步用的是多面体结构：约束条数有限，所以``紧约束多一条''这个计数不能无限进行。第三步用的是不含直线：否则可以沿 \(d\) 永远走下去，永远碰不到新约束。全程没有用到可行域有界——无界多面体只要最优值有限，结论照样成立。

顶点数量随维数指数增长，穷举不可行，但这条命题把搜索范围从连续区域压缩成了一张有限的顶点表，并且相邻顶点之间可以沿边行走。单纯形法走的就是这条路：从一个顶点出发，沿目标下降的边挪到相邻顶点，直到所有出边都不再下降。内点法换了个思路，从可行域内部往下走，用障碍函数把自己挡在边界以内，一路逼近最优点却不踩顶点，做法见\hyperref[sec:09-Convex-Optimization/06-Optimization-Algorithms/05]{``从障碍函数到原始---对偶内点法''}。两条路径殊途同归，都得益于同一个几何事实。

右图那种情形也要留个心眼。等值线与一条边平行时，整条边上的目标值相同，求解器返回哪一端取决于实现细节，换个版本可能换个答案。如果你的下游还在意稀疏性、公平性或者解在数据轻微扰动下的稳定性，这些偏好必须写进模型——加一个次级目标或者一个正则项——而不是指望求解器替你挑。

\subsection{辅助变量把不可微的目标压回线性}

线性规划的表达力远超``几条直线加一个线性目标''的第一印象。流量守恒、供需平衡、配比上下限、运输成本、概率分布的归一化，都是线性等式或不等式。更值得注意的是，一些明显不线性的目标可以被改写进来。

想最小化若干仿射误差中的最大值 \(\max_i(a_i^{\trans}x+b_i)\)。这个函数凸（一族仿射函数的逐点最大值，理由见\hyperref[sec:09-Convex-Optimization/03-Convex-Functions/01]{``弦在图像上方意味着什么''}），但它在几条仿射函数交汇处不可微，求解器也没有直接吃这种表达式的接口。引入一个新变量 \(t\)，要求它压住每一项：

\[
\begin{aligned}
\text{minimize}\quad &t\\
\text{subject to}\quad &a_i^{\trans}x+b_i\le t,\quad i=1,\ldots,m.
\end{aligned}
\]

对固定的 \(x\)，满足全部不等式的最小 \(t\) 恰是那些仿射函数的最大值，所以在 \((x,t)\) 上最小化 \(t\) 与在 \(x\) 上最小化最大值给出同一个 \(x\)。上境图变换的实用形态就是这个：把函数值抬成一个变量，把``取最大''摊成一组不等式。

绝对值走同一条路。\(\abs{a^{\trans}x-b}\le t\) 等价于 \(a^{\trans}x-b\le t\) 与 \(b-a^{\trans}x\le t\) 同时成立。于是 \(\ell_1\) 拟合与 \(\ell_\infty\) 最坏误差都能写成线性规划，多出来的只是几个辅助变量和几行约束。这类改写不是数学游戏，它决定了一个模型能不能进入现成工具链——下一篇会看到，选 \(\ell_1\) 还是 \(\ell_2\) 是一个建模判断，而这里的改写保证了这个判断不会被求解能力否决。

\subsection{求解器返回的不只是一个点}

回到开头那份排产表。写成线性规划以后，你拿到的除了一个方案，还有一组对偶变量：每条约束配一个乘子，量化``这条约束松一点，成本能降多少''。工时上限对应的乘子就是这条产线的边际价值，为零说明这条线还有余量，为正说明它是当前的瓶颈。

更要紧的是，这组乘子构成一份可以独立核验的\textbf{证书}。它给出一个对所有可行点都成立的下界，而线性规划里这个下界恰好等于最优值。接手的人不需要相信你的代码、你的迭代次数或者你的收敛判据，只要把乘子代进几个不等式核对一遍，就能确认这个方案没有改进空间。下界从哪里来、什么时候能与最优值闭合，见\hyperref[sec:09-Convex-Optimization/05-Duality/01]{``Lagrangian 怎样制造全局下界''}与\hyperref[sec:09-Convex-Optimization/05-Duality/02]{``对偶间隙何时闭合''}；乘子作为边际价格的读法见\hyperref[sec:09-Convex-Optimization/05-Duality/04]{``对偶变量是约束的影子价格''}。

线性规划把凸优化的全部要件摆在了最透明的位置：可行域是多面体，目标是平面族，最优解在顶点，对偶是边际价格。后面几篇要做的是把这些几何对象逐步换掉——多面体换成椭球，平面族换成碗，再往后换成半正定锥——而分离、对偶、证书这套逻辑一直没变。

换掉的第一步，是把线性目标改成平方和。那是数据工作者最熟悉的一个式子。

